% !TEX program = xelatex
\documentclass[12pt,a4paper]{ctexart}

\usepackage{fontspec}
\setCJKmainfont{SimSun}[BoldFont=SimHei,ItalicFont=KaiTi]
\setCJKsansfont{SimHei}
\setCJKmonofont{FangSong}

\usepackage[top=28mm,bottom=28mm,left=28mm,right=28mm]{geometry}
\usepackage{setspace}
\usepackage{fancyhdr}
\usepackage{hyperref}
\usepackage{perpage}
\usepackage{xcolor}

\MakePerPage{footnote}

% Footnote numbers in circles
\newcommand{\circledfootnotenumber}{%
  {\fontspec{SimSun}\ifcase\value{footnote}\or ①\or ②\or ③\or ④\or ⑤\or ⑥\or ⑦\or ⑧\or ⑨\or ⑩\else\arabic{footnote}\fi}}
\renewcommand{\thefootnote}{\circledfootnotenumber}

\hypersetup{colorlinks=true, linkcolor=black, urlcolor=blue!50!black}

\setlength{\parindent}{2em}
\setlength{\parskip}{0pt}
\linespread{1.35}

\pagestyle{fancy}
\fancyhf{}
\fancyfoot[C]{\thepage}
\renewcommand{\headrulewidth}{0pt}

\begin{document}

\begin{center}
    {\zihao{2}\bfseries 从《摆脱贫困》看以人民为中心的发展思想\par}
    \vspace{0.5em}
    {\zihao{3}\bfseries ——兼论精准扶贫的基层治理逻辑\par}
    \vspace{1.2em}
    {\zihao{-4}
    姓名：\underline{\makebox[3cm]{纪文龙}} \qquad
    学号：\underline{\makebox[4cm]{2024012842}} \qquad
    班级：\underline{\makebox[3cm]{行健烽火4}}\par}
    \vspace{1.0em}
\end{center}

\begin{abstract}
《摆脱贫困》记录闽东地区走出贫困的历程，是理解中国特色社会主义现代化基层治理的生动实践。它将宏大叙事与群众路线、干部作风建设、调查研究方法以及地方产业选择等内容融为一体。结合课程的相关教学内容，本文聚焦于“以人民为中心”的实践落地，分析《摆脱贫困》中的基层治理逻辑。
通过对脱贫攻坚实践的分析，本文认为：中国脱贫成就的密码，并不在于单纯的行政资源投入，而在于建立了一套响应群众真实需求、广泛动员群众参与、良性互动的制度化和系统化的治理机制。
\end{abstract}

\noindent \textbf{关键词：}《摆脱贫困》；以人民为中心；脱贫攻坚；基层治理；共同富裕

\vspace{1.5em}

\section*{一、问题的提出：为什么从“摆脱贫困”理解以人民为中心}

在学习“习近平新时代中国特色社会主义思想概论”这门课时，我脑子里经常会冒出一个问号：“坚持以人民为中心”这个听上去很大、很宏观的命题，在实际的基层治理中，到底是怎么落地的？在传统的体制或发展经济学视野中，政策往往偏好总量数据与抽象的效率指标。老实说，在没读《摆脱贫困》前，我习惯于从新闻或发展经济学的宏观指标来看扶贫。但仔细想想，冰冷的总量数据和实际的人间冷暖其实隔着很远的距离。

当视角拉回到真实的闽东大山，或者中西部那些山沟沟里时，老百姓需要的并不是宏观报表，而是柴米油盐的生计。比如，水井能不能出水？娃娃上学要走多久？家里种的茶叶能不能卖出去？这些看似琐碎的细节，才是一个家庭最真实的生存处境。

教材中指出，人民立场是新时代中国特色社会主义思想的根本立场，坚持以人民为中心是这一思想的根本出发点和落脚点。这告诉我们，评估任何公共政策的成效，都需要建立在微观个体的生活福祉与社会的客观改善之上，而非仅停留在报表上的数字增长。

正基于此，我选择将《摆脱贫困》这本书作为分析起点。它是一份生动的文献，展现了地方领导者在错综复杂的人文环境、恶劣的自然条件下，如何开展具体的行政管理、政策制定与群众工作。在基层实际运作中，贫困绝非只有收入这单一维度，它是由地理屏障、基础设施缺失、传统观念、基层组织涣散等多重因素交织而成的复杂现状。应对这种深层次困境，“发钱发物”式的救济或行政命令式的摊派，往往会陷入困境。只有将资源的投入、基层党组织的重构以及农户个体能动性结合，才能在泥泞中开辟出可持续发展的道路。

这就使闽东的脱贫实践体现了“人民”价值内核。在这一治理场景中，群众逐渐从被动接受行政资源的“客体”转化为具备内生发展动力的“主体”；与此相应，基层干部也需要将工作重心下沉，在田间地头扮演好协调社会关系、提供公共服务的服务者角色。这种治理角色的调适，揭示了新时代基层治理的规律。

\section*{二、《摆脱贫困》中的价值线索：为民造福不是空泛口号}

通读整本书，我感触最深的一条主线就是公权力与人民福利之间的实际联结。以前看“为官一场，造福一方”这类话，总觉得有点像官样文章或者传统的道德说教，但读进去之后才发现，这其实是在谈一种现代公共行政的常识：当官的手里握着权力，就必须得服务于老百姓，把权力转化为群众生活实实在在的改善。\footnote{习近平：《摆脱贫困》，福建人民出版社，1992年。}这种底层逻辑契合课程中人民主体的思想内涵。从微观的治理实践来看，我们能够深刻体会到“人民”并不是一个抽象的宏大概念，而是具体到每一个对吃水、道路改造、下一代入学以及农产品销路有所期盼的普通农户。

在地方政府的具体实操中，这种“造福”其实很考验干部的作风。我把它粗略归纳为两点：首先是得戒掉形式主义的浮躁。有些干部在面对贫困时，容易觉得是“群众脑子太封建、思想保守”或者干脆怪“穷山恶水没办法”，最后索性搞点刷墙铺路的形象工程来敷衍。其次是要防止“包办代替”。如果只是一味地送钱送物资，而不去扶持村民的主体产业，村民自己没有挣钱的本事，那政府资金一旦退场，大家还是会掉回贫困线以下。学术界在梳理宁德时期的治理模式时，也常将其提炼为一种以人民根本利益为核心的治理模式。\footnote{陈振明、吕志奎：《〈摆脱贫困〉中的地方治理思想研究》，《马克思主义与现实》2017年第1期。}它要求干部将价值关怀内化为日常的行动，在制度运行中体现对人的尊重。

这也引发了我们对公共治理成效的反思：我们很容易只关注资金与项目的单向输入，却容易忽视群众的可持续自我发展能力。如果一个贫困村庄虽然账面上收入达标，但在产业基础、乡村教育、产业协作组织方面依然是一片荒漠，那么这种脱贫就是无源之水、无本之木，难以为继。由此，摆脱绝对贫困与实现共同富裕是两个相互联系的递进进程：前者是通过底线兜底来解决最基本的生存与温饱困难，后者则是通过优化资源配置与提升发展能力，让每个人都能平等地参与和分享现代化的成果。

\section*{三、精准扶贫的治理逻辑：把“人民”具体到每一户、每一项需求}

说实话，我国从早期的“大水漫灌”扶贫转向后来的“精准滴灌”，在治理逻辑上是一个非常大的坎。以前那种修路建桥的开发性扶贫虽然有好处，但由于不够精细，钱花出去了，山沟沟里的那些特困户却可能根本没享受到。学者汪三贵也指出过，精准扶贫在基层落地，最难的关卡就是怎么识别、怎么帮扶、怎么考核，每一个环节都得跟基层治理的实际结合起来。\footnote{汪三贵、郭子豪：《论中国的精准扶贫》，《贵州社会科学》2015年第5期。}这与概论课中强调的“实事求是”和“问题导向”方法论完全契合。如果无法建立起一套深入基层的真实信息网络，政策资源就很容易在层层传导中被耗散，造成资源的浪费。

精准扶贫最了不起的地方，其实就是它打破了宏观的“平均数”。比如一个村子平均收入达标了，不代表里面每户人都过得下去。它把视线真正投向了具体的个人、家庭和活生生的人。每家穷的原因千差万别：有的是因为家里顶梁柱病倒了，有的是因为住在悬崖边上根本没路，还有的是因为没技术也没劳动力。这种碎片化、差异化的治理诉求，必须通过网格化的“一户一策”来予以回应。

当然，精准扶贫政策在基层实际运行中也伴随着各种现实挑战。例如，葛志军、邢成举等学者在针对西部两村的田野调查中发现，在这一精密行政运作的背后，也存在着贫困户参与权虚化、政策柔性不足、以及基层工作强度过重等问题。\footnote{葛志军、邢成举：《精准扶贫：内涵、实践困境及其原因阐释——基于宁夏银川两个村庄的调查》，《贵州社会科学》2015年第5期。}这些实证材料恰恰警示我们，单凭顶层设计蓝图的完美是不够的，基层执行中的调适能力和对当地实情的反馈机制才是政策成败的关键。

同时，精准扶贫并不是没有问题的理想过程。邓维杰在讨论精准扶贫难点时提到，实践中可能出现规模排斥、区域排斥和识别排斥等问题。\footnote{邓维杰：《精准扶贫的难点、对策与路径选择》，《农村经济》2014年第6期。}这表明，以人民为中心从来都不是一蹴而就的终极状态，而是一个不断发现偏差、校正偏轨的动态纠偏过程。如果底层的建档立卡程序脱离了村民委员会的民主监督与日常公示，变成了干部关门评定的“密室决策”，那么所谓的“精准”就只是数据层面的自娱自乐。越是面临复杂的利益分配，越要依赖信息的公开透明与群众的深度参与。

总体来看，精准扶贫的实质是通过求真务实的逻辑去优化基层治理的行为模式。其核心在于通过深入的田野调研破除拍脑门决策，并在执行中根据农户实际致贫原因分类施策，最终把评估权交给群众，以群众的获得感来衡量政策成效。这种模式的生命力在于将国家行政资源的宏观调配与乡村社会的微观生态深度融合，避免陷入无休止的“填表迎检”而空耗治理效能。

\section*{四、调查研究与群众路线：基层治理不能只在材料里完成}

这就自然扯到了另一个很接地气的话题：调查研究。在宁德工作时，习近平同志带头搞了“四下基层”（信访接待下基层、现场办公下基层、调查研究下基层、政策宣传下基层），这在今天看来依然是打通官僚体制阻碍、联系老百姓的绝佳方法。\footnote{申坤：《新时代背景下习近平调查研究重要论述的价值意义》，《山东干部函授大学学报》2020年第3期。}做调研绝对不能像蜻蜓点水那样搞“走基层”，或者把大部分心思花在写报告、填表这种文字游戏上。真正的群众路线，要求干部得有穿透行政文书、在现场面对真问题的干劲。

如果没了扎实的调研，基层治理就容易走歪。一方面，上面的领导如果摸不清底下的真实温度，光在办公室里看报表、做推演，拿出来的政策多半是脱离实际的；另一方面，底下的干部在各种层层加码的考核下，为了应付检查，不得不把时间和精力都花在写材料、做PPT上，成了所谓的“纸面扶贫”。一户人家真实的难处，哪能仅靠表格里的几个收入数字看出来？一个村子能不能发展，又哪是几张规划图能画出来的？干部必须得坐到老百姓的炕头上，在面对面的聊天里听出那些被官方文件过滤掉的真实声音。

这也就是为什么《摆脱贫困》这本书到今天读起来依然让人热血沸腾。书里没有高高在上的说教，也没有把闽东当成一个需要被施舍的对象，而是实打实地反思干部的作风、梳理老百姓的心态，因地制宜地找路子。现在我们虽然消除了绝对贫困，开始搞“乡村振兴”，但这一套务实的逻辑依然有用。摘帽脱贫可不是“毕其功于一役”，怎么防返贫、怎么把村子里的产业真正盘活，这些都需要更细、更稳的后续措施。

\section*{五、从脱贫攻坚到共同富裕：发展成果怎样更稳定地留在人民生活中}

从大历史的视角来看，脱贫攻坚确实彻底改变了中国农村落后的旧样貌。杨明伟老师曾指出，摆脱贫困并迈向共同富裕，是老百姓最朴素的强国期盼，也是一场了不起的社会变革。\footnote{杨明伟：《百年奋斗史中的摆脱贫困迈向共同富裕》，《红旗文稿》2021年第6期。}然而，站在“中国式现代化”的历史新起点上，脱贫攻坚的胜利并非终点，而是全面推进乡村振兴、迈向共同富裕的起点。共同富裕被赋予了超越传统资本主义福利国家的全新社会理想，它不仅要求财富总量做大，更要求我们在地区之间、阶层之间打通公平竞争和自我发展的通道，给大家提供平等改变命运的机会。

我特别赞同一句话：共同富裕绝不是搞平均主义的“大锅饭”，而是一个一步一个脚印的长期过程。不过，巩固脱贫成果仍然不容易。中西部很多刚摘帽的地方产业单一、家底子薄，教育和医疗资源跟大城市比还是有不小的差距。研究者胡蝶也提过，只有把教育资源真正精准地配到刀刃上，才能彻底掐断贫困在代际之间的恶性循环，给困难群众自我发展的底气。\footnote{胡蝶：《扶贫必扶智：教育精准扶贫是摆脱贫困的内生动力》，《改革与开放》2016年第21期。}这也说明，真正稳定的脱贫不能只盯着今年账面上过了贫困线没有，关键得看下一代能不能通过公平的教育和健康保障，实现往上走的良性流动。

如果我们把《摆脱贫困》、精准扶贫、乡村振兴与共同富裕置于同一个大历史视角下审视，一条“从生存权保障到发展权共享”的演进脉络便清晰可见。其首要目标是保障基本温饱和生存安全，随后通过教育帮扶与产业导入来提升群众发展能力，最终指向以更加完善的分配制度达成现代化成果的全社会共享。习近平总书记在面对突发公共卫生事件与经济波动等复杂局面时，曾坚决部署必须坚守如期打赢脱贫攻坚战的底线目标，深刻揭示了将人民生活水平与社会正义置于最优先地位的执政理念。\footnote{习近平：《在决战决胜脱贫攻坚座谈会上的讲话》（2020年3月6日），《求是》2020年第6期。}

作为一名大学生，在思考这些治理逻辑时，我觉得最应该防范的就是空谈概念、自说自话。我们得把眼光投向身边真实的细微变化。比如：回乡卖农产品的师兄师姐，到底有没有拿到政府的政策补贴？村里小学撤并之后，住在远山里的孩子上学是不是更难了？镇上的年轻干部天天写材料迎检，还能有多少时间和心思真正去村里访贫问苦？这些不起眼的小事，才是衡量“以人民为中心”有没有在泥土里扎根的试金石。

\section*{结语}

尽管书里的一些具体数字和年代背景已经离我们有点远了，但《摆脱贫困》给我的启发依然很深刻。它证明了“以人民为中心”绝不仅仅是一句挂在墙上的口号或政治修辞，它实实在在地体现在干部下基层走访的脚步里、体制机制的自我纠错里，还有跟老百姓心贴心的日常交流中。扶贫这件大事，无非就是把抽象的“人民”二字，细化为每一个被生活困难困扰的家庭和具体的人。

任何大政策在基层落地，难免都会遇到点摩擦或体制上的卡壳。要想解决这些问题，关键还是在于能不能实事求是，敢不敢直面老百姓的各种急难愁盼。从早期的闽东脱贫，到今天的乡村振兴，这条线其实一直都没变：就是让现代化建设的好处，能够顺畅地润泽进每一个普通人的日常生活里。

\end{document}
